\chapterAndLabel{Modifiability Checker}{modifiability-checker}

\newcommand{\UnsupportedOperationException}{\<Un\-sup\-port\-ed\-Op\-er\-a\-tion\-Ex\-cep\-tion>}

In Java, an \emph{optional} method is one that a class may or may not
support.  (Optional methods are unrelated to the \<Optional> type.)
An example is
\sunjavadoc{java.base/java/util/Collection.html\#add(E)}{add}.  Some
subclasses of \<Collection> support \<add>, and others (such as the
result of calling
\sunjavadoc{java.base/java/util/Collections.html\#unmodifiableCollection(java.util.Collection)}{Collections.unmodifiableCollection})
do not support \<add>.

Calling an unsupported method throws \UnsupportedOperationException.  Thus,
whenever a client calls an optional method, there is a possibility of an
exception.  The Modifiability Checker eliminates that possibility.

To run the Modifiability Checker:

\begin{Verbatim}
javac -processor org.checkerframework.checker.modifiability.ModifiabilityChecker MyFile.java
\end{Verbatim}

If the Modifiability Checker issues no warning, then your program is
guaranteed not to throw \UnsupportedOperationException\ due to invoking an
optional method (such as
\sunjavadoc{java.base/java/util/Collection.html\#add(E)}{add} or
\sunjavadoc{java.base/java/util/Collection.html\#remove(java.lang.Object)}{remove})
on an unmodifiable collection.
(For brevity, this chapter generally uses
just ``collection'' to mean ``collection, map, or iterator''.)

It is possible to write a checker that warns about other unsupported
methods (beyond those in the Java Collections Framework), but the
Modifiability Checker does not do so.


\sectionAndLabel{Unmodifiable collections}{modifiability-unmodifiable-collections}

An \emph{unmodifiable collection} is one that throws
\UnsupportedOperationException\ if a mutating method is called.
The JDK provides many such collections, for example, the results of:
\begin{itemize}
\item
  \sunjavadoc{java.base/java/util/List.html\#of(E...)}{\code{List.of()}},
  \sunjavadoc{java.base/java/util/List.html\#copyOf(java.util.Collection)}{\code{List.copyOf()}}
\item
  \sunjavadoc{java.base/java/util/Set.html\#of(E...)}{\code{Set.of()}},
  \sunjavadoc{java.base/java/util/Set.html\#copyOf(java.util.Collection)}{\code{Set.copyOf()}}
\item
  \sunjavadoc{java.base/java/util/Map.html\#of(K,V,K,V)}{\code{Map.of()}},
  \sunjavadoc{java.base/java/util/Map.html\#copyOf(java.util.Map)}{\code{Map.copyOf()}}
\item
  \sunjavadoc{java/util/Collections.html\#emptyList()}{\code{Collections.emptyList()}}
  \sunjavadoc{java/util/Collections.html\#emptySet()}{\code{Collections.emptySet()}}
  \sunjavadoc{java/util/Collections.html\#emptyMap()}{\code{Collections.emptyMap()}}
\item
  \sunjavadoc{java.base/java/util/Collections.html\#unmodifiableList(java.util.List)}{\code{Collections.unmodifiableList()}}
  \sunjavadoc{java.base/java/util/Collections.html\#unmodifiableSet(java.util.Set)}{\code{Collections.unmodifiableSet()}}
  \sunjavadoc{java.base/java/util/Collections.html\#unmodifiableMap(java.util.Map)}{\code{Collections.unmodifiableMap()}}
\end{itemize}

A \emph{modifiable collection} is one on which mutating methods can be called
without throwing \UnsupportedOperationException.  Non-optional methods, such
as \code{size}, \code{get}, \code{contains}, \code{isEmpty}, and \code{toArray}, may be
called on any collection, whether modifiable or unmodifiable.

Some collections may be modified in some ways but not in others.
A collection can be modified in three different ways:
\textbf{Grow}, \textbf{Shrink}, and \textbf{Replace}.
\begin{itemize}
    \item \textbf{Grow}: The collection can grow by adding new elements
      (\<add>, \<addAll>, \<offer>, \ldots).
    \item \textbf{GrowSequenced}: The collection can grow by adding new
      elements at the beginning or end (\<addFirst>, \<addLast>).
      Because GrowSequenced checking is exactly like Grow checking for a
      different set of methods, this chapter often omits it for brevity.
    \item \textbf{Shrink}: The collection can shrink by removing
      existing elements (\<remove>, \<removeAll>, \<clear>, \<retainAll>, \<poll>, \<pop>, \ldots).
    \item \textbf{Replace}: The collection can replace existing
      elements with new values at specific locations/keys (\<List.set>, \<Map.replace>, \<Map.replaceAll>,
       \<ListIterator.set>,  \<Collections.sort>, \<Map.Entry.setValue>, \ldots).
\end{itemize}

\noindent
Some collections can also be modified through their iterator; see Section~\ref{modifiability-iterator-behavior}.



\sectionAndLabel{Modifiability annotations}{modifiability-annotations}

The Modifiability Checker uses \textbf{multiple independent 4-element
hierarchies}, one per capability.  Every collection type carries exactly
one qualifier from each hierarchy.  The checker also uses an additional
2-element hierarchy to express whether \<remove()> can be called on
a collection's iterator.

% TODO: add a figure that follows the sturucture of other checker manuals.
\begin{figure}
\begin{smaller}
  \centering
  \begin{minipage}{0.85\linewidth}
\begin{verbatim}
  Grow hierarchy           Shrink hierarchy              Replace hierarchy             Iterator hierarchy

    @MaybeGrowable             @MaybeShrinkable               @MaybeReplaceable          @MaybeIteratorPolyMod
      /      \                 /       \                    /       \                          |
 @Growable @Ungrowable   @Shrinkable @Unshrinkable   @Replaceable @Unreplaceable               |
      \      /                 \       /                    \       /                          |
    @BottomGrowable             @BottomShrinkable                @BottomReplaceable              @IteratorPolyMod
\end{verbatim}
  \end{minipage}
\end{smaller}
  \caption{Independent modifiability qualifier hierarchies.
  (The GrowSequenced hierarchy is omitted for brevity.)
    Each annotated type carries one qualifier from each hierarchy.}
  \label{fig-modifiability-lattice}
\end{figure}

Figure~\ref{fig-modifiability-lattice} shows the hierarchies.
Each hierarchy is type-checked individually.
An assignment \code{T x = y} is valid only if the compile-time type of \code{y} is a
subtype of (or equal to) \code{T} in every hierarchy simultaneously.  For
example, \code{\emph{@Growable} @MaybeShrinkable @MaybeReplaceable List} is a
subtype of \code{\emph{@MaybeGrowable} @MaybeShrinkable @MaybeReplaceable List}
(because \<@Growable> $\le$ \<@MaybeGrowable>), but it is \emph{not} a
subtype of \code{\emph{@Ungrowable} @MaybeShrinkable @MaybeReplaceable List}
(because \<@Growable> and \<@Ungrowable> are incomparable siblings in the
Grow hierarchy).

\subsectionAndLabel{Alias annotations}{modifiability-aliases}

For convenience, the Modifiability Checker provides four alias
annotations.  They are not part of any hierarchy; the checker expands
each one into its constituent qualifiers when it appears in source code.

\begin{description}

\item[\refqualclass{checker/modifiability/qual}{Modifiable}]
Alias for \code{@Growable @Shrinkable @Replaceable} (but see
Section~\ref{modifiability-types-lacking-methods} for exceptions).  Calling
any mutating operation on this collection will not throw
\UnsupportedOperationException.

\item[\refqualclass{checker/modifiability/qual}{Unmodifiable}]
Alias for \code{@Ungrowable @Unshrinkable @Unreplaceable} (but see
Section~\ref{modifiability-types-lacking-methods} for exceptions).  The
collection is definitely unmodifiable: calling any mutating operation
always throws \UnsupportedOperationException.  The checker issues a warning
at every invocation of a grow, shrink, or replace operation.

\item[\refqualclass{checker/modifiability/qual}{MaybeModifiable}]
Alias for \code{@MaybeGrowable @MaybeShrinkable @MaybeReplaceable}.
The checker cannot determine the modifiability of the collection.
The checker conservatively issues a warning at every invocation of a grow,
shrink, or replace operation.

\item[\refqualclass{checker/modifiability/qual}{UnmodifiableParam}]
Syntactic sugar for \refqualclass{checker/modifiability/qual}{MaybeModifiable}.
It may only be written within a formal parameter type to indicate that the
method does not modify the parameter.  This annotation (and
\refqualclass{checker/modifiability/qual}{MaybeModifiable}) enable the
programmer to distinguish between a data structure's capability to perform
a modification and the programmer's intent about whether a modification
should occur.
\refqualclass{checker/modifiability/qual}{UnmodifiableParam}
may appear nested within a parameter type, such as in
\code{List<@UnmodifiableParam List<String>\relax>}.

\refqualclass{checker/modifiability/qual}{UnmodifiableParam} represents
\emph{reference unmodifiability}:  no modification may happen through the
reference, but an alias may modify the value.  By contrast,
\refqualclass{checker/modifiability/qual}{Unmodifiable} is object
unmodifiability:  no alias may modify the value.

\item[\refqualclass{checker/modifiability/qual}{PolyModifiable}]
Alias for \code{@PolyGrowable @PolyShrinkable @PolyReplaceable}
(see Section~\ref{modifiability-polymorphic}).

\end{description}

\noindent
Examples using alias annotations:

\begin{Verbatim}
  @Modifiable List<String> mod = new ArrayList<>();  // all three capabilities
  @Unmodifiable List<String> unmod = List.of("a");   // no capabilities
  @MaybeModifiable List<String> unknown = ...;  // capabilities are unknown at compile time
  void foo(@UnmodifiableParam List<String> list) { ... }  // foo does not modify list
\end{Verbatim}

\subsectionAndLabel{Qualifiers}{modifiability-qualifiers}

\subsubsectionAndLabel{The Grow hierarchy}{modifiability-grow-hierarchy}

\begin{description}
\item[\refqualclass{checker/modifiability/qual}{MaybeGrowable}]
The top qualifier in the Grow hierarchy.
The checker cannot determine whether this collection supports grow
operations.
Calling grow operations such as \code{add}, \code{addAll}, etc.\ on this
collection \emph{may} throw \UnsupportedOperationException.
This is the default qualifier for unannotated types in the
Grow hierarchy.

\item[\refqualclass{checker/modifiability/qual}{Growable}]
Calling grow operations such as \code{add}, \code{addAll}, etc.\ on this
collection \emph{will not} throw \UnsupportedOperationException.

\item[\refqualclass{checker/modifiability/qual}{Ungrowable}]
Calling grow operations such as \code{add}, \code{addAll}, etc.\ on this
collection \emph{will} throw \UnsupportedOperationException.

\item[\refqualclass{checker/modifiability/qual}{BottomGrowable}]
The bottom qualifier in the Grow hierarchy.  Programmers should rarely write it.
\end{description}

\subsubsectionAndLabel{The GrowSequenced hierarchy}{modifiability-grow-hierarchy}

\begin{description}
\item[\refqualclass{checker/modifiability/qual}{MaybeGrowableSequenced}]
The top qualifier in the GrowSequenced hierarchy.
The checker cannot determine whether this collection supports grow
operations.
Calling sequenced grow operations such as \code{addFirst} and \code{addLast} on this
collection \emph{may} throw \UnsupportedOperationException.
This is the default qualifier for unannotated types in the
GrowSequenced hierarchy.

\item[\refqualclass{checker/modifiability/qual}{GrowableSequenced}]
Calling sequenced grow operations such as \code{addFirst} and \code{addLast} on this
collection \emph{will not} throw \UnsupportedOperationException.

\item[\refqualclass{checker/modifiability/qual}{UngrowableSequenced}]
Calling sequenced grow operations such as \code{addFirst} and \code{addLast} on this
collection \emph{will} throw \UnsupportedOperationException.

\item[\refqualclass{checker/modifiability/qual}{BottomGrowableSequenced}]
The bottom qualifier in the GrowSequenced hierarchy.  Programmers should rarely write it.
\end{description}

\subsubsectionAndLabel{The Shrink hierarchy}{modifiability-shrink-hierarchy}

\begin{description}
\item[\refqualclass{checker/modifiability/qual}{MaybeShrinkable}]
The top qualifier in the Shrink hierarchy.
The checker cannot determine whether this collection supports shrink
operations.
Calling shrink operations such as \code{remove}, \code{clear}, etc.\ on
this collection \emph{may} throw \UnsupportedOperationException.
This is the default qualifier for unannotated types in the
Shrink hierarchy.

\item[\refqualclass{checker/modifiability/qual}{Shrinkable}]
Calling shrink operations such as \code{remove}, \code{clear}, etc.\ on
this collection \emph{will not} throw \UnsupportedOperationException.

\item[\refqualclass{checker/modifiability/qual}{Unshrinkable}]
Calling shrink operations such as \code{remove}, \code{clear}, etc.\ on
this collection \emph{will} throw \UnsupportedOperationException.

\item[\refqualclass{checker/modifiability/qual}{BottomShrinkable}]
The bottom qualifier in the Shrink hierarchy.  Programmers should rarely write it.
\end{description}

\subsubsectionAndLabel{The Replace hierarchy}{modifiability-replace-hierarchy}

\begin{description}
\item[\refqualclass{checker/modifiability/qual}{MaybeReplaceable}]
The top qualifier in the Replace hierarchy.
The checker cannot determine whether this collection supports replace
operations.
Calling replace operations such as \code{set}, \code{replaceAll}, etc.\
on this collection \emph{may} throw \UnsupportedOperationException.
This is the default qualifier for unannotated types in the
Replace hierarchy.

\item[\refqualclass{checker/modifiability/qual}{Replaceable}]
Calling replace operations such as \code{set}, \code{replaceAll}, etc.\
on this collection \emph{will not} throw \UnsupportedOperationException.

\item[\refqualclass{checker/modifiability/qual}{Unreplaceable}]
Calling replace operations such as \code{set}, \code{replaceAll}, etc.\
on this collection \emph{will} throw \UnsupportedOperationException.

\item[\refqualclass{checker/modifiability/qual}{BottomReplaceable}]
The bottom qualifier in the Replace hierarchy.  Programmers should rarely write it.
\end{description}

\subsubsectionAndLabel{The Iterator hierarchy}{modifiability-iterator-hierarchy}

\begin{description}
\item[\refqualclass{checker/modifiability/qual}{MaybeIteratorPolyMod}]
The top qualifier.
The checker cannot determine whether this collection's \code{iterator()} is
\<@PolyShrinkable>.
This is the default qualifier for unannotated types.

\item[\refqualclass{checker/modifiability/qual}{IteratorPolyMod}]
This collection's \code{iterator()} is \<@PolyShrinkable>.  That is, if
collection \<c> is \<@Shrinkable>, then \<c.iterator()> is also
\<@Shrinkable>.

\end{description}


\subsectionAndLabel{Polymorphic qualifiers}{modifiability-polymorphic}

A polymorphic qualifier (Section~\ref{method-qualifier-polymorphism})
specifies that a method's return type has the
same qualifier (in some hierarchy) as a formal parameter type (possibly
the receiver).

\begin{description}

\item[\refqualclass{checker/modifiability/qual}{PolyGrowable}]
Polymorphic qualifier for the Grow hierarchy.
The return type's growability matches the growability of whichever
formal parameter is annotated with \<@PolyGrowable>.

\item[\refqualclass{checker/modifiability/qual}{PolyShrinkable}]
Polymorphic qualifier for the Shrink hierarchy.
The return type's Shrink qualifier matches the Shrink qualifier of
whichever formal parameter is annotated with \<@PolyShrinkable>.
Useful for methods such as \code{Map.keySet()} that preserve shrinkability while
returning an \<@Ungrowable> view.

\item[\refqualclass{checker/modifiability/qual}{PolyReplaceable}]
Polymorphic qualifier for the Replace hierarchy.
The return type's replaceability matches the replaceability of
whichever formal parameter is annotated with \<@PolyReplaceable>.

\item[\refqualclass{checker/modifiability/qual}{PolyModifiable}]
Alias for \code{@PolyGrowable @PolyShrinkable @PolyReplaceable}.
Preserves all three capabilities simultaneously.
Use on methods such as \code{Collections.synchronizedList} thatx
do not change modifiability.

\end{description}


\subsectionAndLabel{Modifiability method annotations}{modifiability-method-annotations}

The Modifiability Checker supports one annotation that specifies method
behavior.

\begin{description}

\item[\refqualclass{checker/modifiability/qual}{ThrowsUOE}]
When written on a
method or constructor declaration, it indicates that the method always throws
\UnsupportedOperationException.  Therefore, any code that calls it is
erroneous.  The Modifiability Checker issues an error at every call to a
method annotated with \code{@ThrowsUOE}.

\end{description}

\sectionAndLabel{\<@Modifiable> and \<@Unmodifiable> for types without grow, shrink, and/or replace methods}{modifiability-types-lacking-methods}

%% TODO add explicit examples TODO, also in the Javadoc TODO.

Ordinarily, the declaration
\<@Modifiable MyClass x;> is equivalent to
\<@Growable @Shrinkable @Replaceable MyClass x;>,
and
\<@Unodifiable MyClass x;> is equivalent to
\<@Ungrowable @Unshrinkable @Unreplaceable MyClass x;>.  However, there are
exceptions:

%% TODO add more examples, exhaustively if possible.
\begin{tabular}{rlcrrrl}
% \<@Modifiable X>      & = & \<@Growable> & \<@Shrinkable> & \<@Replaceable> & \<X> \\
\<@Modifiable> & \<Queue>    & = & \<@Growable> & \<@Shrinkable> & \<@MaybeReplaceable> & \<Queue> \\
\<@Unmodifiable> & \<Queue>    & = & \<@Ungrowable> & \<@Unshrinkable> & \<@MaybeReplaceable> & \<Queue> \\
\<@Modifiable> & \<Set>      & = & \<@Growable> & \<@Shrinkable> & \<@MaybeReplaceable> & \<Set> \\
\<@Unmodifiable> & \<Set>      & = & \<@Ungrowable> & \<@Unshrinkable> & \<@MaybeReplaceable> & \<Set> \\
\<@Modifiable> & \<Iterator> & = & \<@MaybeGrowable> & \<@Shrinkable> & \<@MaybeReplaceable> & \<Iterator> \\
\<@Unmodifiable> & \<Iterator> & = & \<@MaybeGrowable> & \<@Unshrinkable> & \<@MaybeReplaceable> & \<Iterator> \\
\end{tabular}

When a programmer writes \<@Modifiable> or \<@Unmodifiable> on a type that
has no methods for a given mutation, the Modifiability Checker uses the top
qualifier.
For example, \sunjavadoc{java.base/java/util/Queue.html}{\code{Queue}} has
\<add> and \<remove> methods but no replacement methods.


Although no object of class \<Queue> supports replace operations,
the type
\<@Replaceable Queue> still makes sense, because some objects of type
\<Queue> do support replace operations.
For example, \<@Replaceable Queue> is the type of
\<q> in \<Queue q = new LinkedList();>.
This behavior ensures that code like this type-checks:

\begin{Verbatim}
@Growable LinkedList list = ...;
Queue q = list;
@Growable LinkedList list2 = (Queue) q; // still @Growable
\end{Verbatim}


\sectionAndLabel{Behavior of the \<iterator()> method}{modifiability-iterator-behavior}

By default, the signature of \<Collection.iterator()> is

\begin{Verbatim}
@Unshrinkable Iterator iterator()
\end{Verbatim}

However, for some subclasses of \<Collection>, the signature of
\<iterator()> is

\begin{Verbatim}
@PolyShrinkable Iterator iterator(@PolyShrinkable MySubclassOfCollection this)
\end{Verbatim}

A collection uses the second signature if the collection is
\<@IteratorPolyMod>.  The collection uses the first signature if the
collection is \<@MaybeIteratorPolyMod>.

Here are examples of its use:

\begin{Verbatim}
  @IteratorPolyMod @Shrinkable List<String> list = ...;
  @Shrinkable Iterator<String> it = list.iterator();
  it.remove();  // OK

  @Unshrinkable List<String> list2 = ...;
  @Unshrinkable Iterator<String> it2 = list2.iterator();
  it2.remove();  // Error: it2 is @Unshrinkable
\end{Verbatim}



\sectionAndLabel{Examples}{modifiability-examples}

\begin{Verbatim}
import java.util.*;
import org.checkerframework.checker.modifiability.qual.*;

class Demo {

  void basicUsage() {
    // @Modifiable is an alias for @Growable @Shrinkable @Replaceable
    @Modifiable List<String> mod = new ArrayList<>();
    mod.add("a");     // OK: @Growable guarantees add() works
    mod.remove("a");  // OK: @Shrinkable guarantees remove() works
    mod.set(0, "b");  // OK: @Replaceable guarantees set() works

    // @Unmodifiable is an alias for @Ungrowable @Unshrinkable @Unreplaceable
    @Unmodifiable List<String> unmod = List.of("a", "b");
    unmod.add("c");   // Error: add() requires @Growable, got @Ungrowable
    unmod.get(0);     // OK: read-only access needs no capability
  }

  void fineGrainedPermissions(@Growable List<String> g,
                              @Shrinkable List<String> s,
                              @Replaceable List<String> r) {
    // Growable list allows adding elements
    g.add("a");       // OK
    g.remove("a");    // Error: remove() requires @Shrinkable
    g.set(0, "b");    // Error: set() requires @Replaceable

    // Shrinkable list allows removing elements
    s.remove("a");    // OK
    s.add("a");       // Error: add() requires @Growable

    // Replaceable list allows updating elements
    r.set(0, "b");    // OK
    r.add("b");       // Error: add() requires @Growable
  }

  void combinedPermissions(@Growable @Replaceable Map<String, String> map) {
    // Map.put requires both Grow and Replace capabilities
    map.put("key", "value");  // OK

    // Map.remove requires Shrink capability
    map.remove("key");        // Error: @MaybeShrinkable (default) !<: @Shrinkable
  }

  // @PolyModifiable preserves all three capabilities
  @PolyModifiable List<String> wrap(@PolyModifiable List<String> list) {
    return list;
  }

  void testPoly(@Modifiable List<String> mod, @Growable @Shrinkable List<String> gs) {
    @Modifiable List<String> m = wrap(mod);            // OK
    @Growable @Shrinkable List<String> gs2 = wrap(gs); // OK
    // :: error: [assignment]
    @Modifiable List<String> bad = wrap(gs);  // Error: gs has @MaybeReplaceable
  }
}
\end{Verbatim}

% LocalWords:  Modifiability modifiability copyOf emptyList emptySet gsr
% LocalWords:  emptyMap unmodifiableList unmodifiableSet unmodifiableMap
% LocalWords:  toArray addAll removeAll retainAll replaceAll setValue
% LocalWords:  MaybeModifiable UnmodifiableParam Growable MaybeGrowable MaybeShrinkable
% LocalWords:  MaybeReplaceable growable ListIterator PolyModifiable addFirst
% LocalWords:  synchronizedList PolyShrinkable PolyGrowable PolyReplaceable ThrowsUOE
% LocalWords:  addLast unreplaceable Growability shrinkability replacability
% LocalWords:  growability Shrinkable Replaceable Ungrowable Unshrinkable
% LocalWords:  Unreplaceable BottomGrowable BottomShrinkable BottomReplaceable
% LocalWords:  MaybeIteratorPolyMod IteratorPolyMod
